\input{libraries/report-preamble.tex}
\title{Binary GCD on Unsigned 64-Bit Integers}
\begin{document}
\maketitle

\section{Task and interface}
\texttt{gcdBinary} computes the greatest common divisor of two \texttt{UInt64}
values using shifts, comparisons, and subtraction.  It accepts every pair of
input words.  The intended result is $\gcd(a,b)$, with $\gcd(0,0)=0$.

LeanExe compiles the Lean function and its callers to WebAssembly (WASM).
The examples run that WASM in Wasmtime.

\section{Verification status}
Lean checks the source types.  Execution tests pass for selected inputs:
both GCD implementations agree with an arbitrary-precision integer reference
in WASM and native Lean.  The native test also checks Lean's \texttt{Nat.gcd}.
Formal correctness and termination proofs for this source function,
its generated WASM, and its fraction client remain deferred.

\section{Algorithm and executable source}
Common factors of two can be removed from both operands and restored at the
end.  A factor of two in just one operand can be removed when the other is odd.
Subtracting the smaller odd operand from the larger produces an even difference.
These identities are described in NIST's binary GCD entry~\cite{binary}.

\lstinputlisting{LeanExe/Lib/NumberTheory/BinaryGcd/Basic.lean}

The listing comes from \texttt{Basic.lean}.  Zero inputs return before a shift
loop can encounter zero.  The common shift count is at most 63 for nonzero
inputs.  Ordering the operands before subtraction gives a nonnegative
difference.  The result fits in \texttt{UInt64}.  The loop keeps two operands
and a shift count.

\section{A shared client task}
The example program provides two fraction-reduction functions.  One calls
Euclidean GCD and the other calls binary GCD.  Both use the same result helper,
which rejects a zero denominator and divides the numerator and denominator
by the selected divisor.  The commands run from the repository root after
following the
\href{https://github.com/jsmorph/leanexe/blob/lib1/DEVELOPING.md}{setup guide}.

\begin{lstlisting}
tools/seminum gcd-binary 48 18
6
tools/seminum fraction-binary 48 18
8/3
tools/seminum fraction 48 18
8/3
\end{lstlisting}

Arguments are unsigned decimal words through $2^{64}-1$.  Invalid syntax
and a zero denominator produce status 2 and a diagnostic on standard error.

\section{Tests}
From the repository root, \texttt{node test/seminum.js} runs the comparisons
described above on the same word pairs for both GCD implementations.
Inputs include zero, maximum words, powers of two, consecutive large Fibonacci
numbers, and generated full-width pairs.  Fraction cases check reduction and
rejection of a zero denominator.  CLI tests check results and invalid-input
diagnostics.

The test source is \texttt{test/seminum.js}.  The native reference driver is
\texttt{test/SeminumNative.lean}.

\begin{thebibliography}{9}
\bibitem{binary} Paul E. Black, ``binary GCD,'' NIST Dictionary of Algorithms
and Data Structures, 2020.
\url{https://xlinux.nist.gov/dads/HTML/binaryGCD.html}.
This entry supplies the parity identities and identifies Knuth's Algorithm
4.5.2B.  The implementation here counts shifts and handles zero inputs before
the loop.
\end{thebibliography}
\end{document}
